%!TEX root = ../main.tex

\chapter{Lưu đồ thuật toán}
\hypertarget{localtoc\thechapter}{}
\localtoc    

\begin{problem}
    Tính $S(n) = 1 + 2 + 3 + \cdots + n$ 
\end{problem}

\code{code/ch01/1-0-1.cpp}{c++}

\begin{problem}
    Tính $S(n) = 1^2 + 2^2 + \cdots + n^2$
\end{problem}

\code{code/ch01/1-0-2.cpp}{c++}

\begin{problem}
    Tính $S(n) = 1 + ½ + \frac{1}{3} + \cdots + \frac{1}{n}$
\end{problem}

\code{code/ch01/1-0-3.cpp}{c++}

\begin{problem}
    Tính $S(n) = ½ + ¼ + \cdots + \frac{1}{2n}$
\end{problem}

\code{code/ch01/1-0-4.cpp}{c++}

\begin{problem}
    Tính $S(n) = 1 + \frac{1}{3} + \frac{1}{5} + \cdots + \frac{1}{2n + 1}$
\end{problem}

\code{code/ch01/1-0-5.cpp}{c++}

\begin{problem}
   Tính $S(n) = \frac{1}{1}×2 + \frac{1}{2}×3 +\cdots+ \frac{1}{n}\times(n + 1)$
\end{problem}

\code{code/ch01/1-0-6.cpp}{c++}

\begin{problem}
    Tính $S(n) = \frac{1}{2} + \frac{2}{3} + \frac{3}{4} + \cdots + \frac{n}{n + 1}$
\end{problem}

\code{code/ch01/1-0-7.cpp}{c++}

\begin{problem}
    Tính $S(n) = \frac{1}{2} + \frac{3}{4} + \frac{5}{6} + \cdots + \frac{2n + 1}{2n + 2} $
\end{problem}

\code{code/ch01/1-0-8.cpp}{c++}

\begin{problem}
   Tính $T(n) = 1 \times 2 \times 3 \times \cdots \times N$
\end{problem}

\code{code/ch01/1-0-9.cpp}{c++}

\begin{problem}
   Tính $T(x, n) = x^n$
\end{problem}

\code{code/ch01/1-0-10.cpp}{c++}

\begin{problem}
    Tính $S(n) = 1 + 1.2 + 1.2.3 + \cdots + 1.2.3 \cdots n$
\end{problem}

\code{code/ch01/1-0-11.cpp}{c++}

\begin{problem}
   Tính $S(n,x) = x + x^2 + x^3 + \cdots + x^n$
\end{problem}

\code{code/ch01/1-0-12.cpp}{c++}

\begin{problem}
  Tính $S(n,x) = x^2 + x^4 + \cdots + x^{2n}$
\end{problem}
\code{code/ch01/1-0-13.cpp}{c++}

\begin{problem}
   Tính $S(n,x) = x + x^3 + x^5 + \cdots + x^{2n + 1}$
\end{problem}
\code{code/ch01/1-0-14.cpp}{c++}

\begin{problem}
   Tính $S(n) = \frac{1}{1} + \frac{1}{1+2} + \frac{1}{1+2+3} + \cdots + \frac{1}{1+2+3+ \cdots + n}$
\end{problem}
\code{code/ch01/1-0-15.cpp}{c++}

\begin{problem}
    Tính $S(n,x) = x + \frac{x^2}{1+2} + \frac{x^3}{1+2+3} + \cdots + \frac{x^n}{1+2+3+\cdots+n}$
\end{problem}
\code{code/ch01/1-0-16.cpp}{c++}

\begin{problem}
    Tính $S(n,x) = x + \frac{x^2}{2!} + \frac{x^3}{3!} + \cdots + \frac{x^n}{n!}$
\end{problem}
\code{code/ch01/1-0-17.cpp}{c++}

\begin{problem}
    Tính $S(n,x) = 1 + \frac{x^2}{2!} + \frac{x^4}{4!} + \cdots + \frac{x^{2n}}{(2n)!}$
\end{problem}
\code{code/ch01/1-0-18.cpp}{c++}

\begin{problem}
    Tính $S(n,x) = 1 + x + \frac{x^3}{3!} + \frac{x^5}{5!} + \cdots + \frac{x^{2n+1}}{(2n+1)!}$
\end{problem}
\code{code/ch01/1-0-19.cpp}{c++}

\begin{problem}
    Liệt kê tất cả “ước số” của số nguyên dương n.
\end{problem}
\code{code/ch01/1-0-20.cpp}{c++}
\begin{problem}
    Tính tổng tất cả “ước số” của số nguyên dương n.
\end{problem}
\code{code/ch01/1-0-21.cpp}{c++}
\begin{problem}
    Tính tích tất cả “ước số” của số nguyên dương n.
\end{problem}
\code{code/ch01/1-0-22.cpp}{c++}
\begin{problem}
    Đếm số lượng “ước số” của số nguyên dương n.
\end{problem}
\code{code/ch01/1-0-23.cpp}{c++}

\begin{problem}
    Liệt kê tất cả “ước số lẻ” của số nguyên dương n.
\end{problem}
\code{code/ch01/1-0-24.cpp}{c++}

\begin{problem}
    Tính tổng tất cả “ước số chẵn” của số nguyên dương n.
\end{problem}
\code{code/ch01/1-0-25.cpp}{c++}

\begin{problem}
    Tính tích tất cả “ước số lẻ” của số nguyên dương n.
\end{problem}
\code{code/ch01/1-0-26.cpp}{c++}

\begin{problem}
    Đếm số lượng “ước số chẵn” của số nguyên dương n.
\end{problem}
\code{code/ch01/1-0-27.cpp}{c++}

\begin{problem}
    Cho số nguyên dương n. Tính tổng các ước số nhỏ hơn chính nó.
\end{problem}
\code{code/ch01/1-0-28.cpp}{c++}

\begin{problem}
    Tìm ước số lẻ lớn nhất của số nguyên dương n. Ví dụ n = 100 ước lẻ
lớn nhất của 100 là 25.
\end{problem}
\code{code/ch01/1-0-29.cpp}{c++}

\begin{problem}
    Cho số nguyên dương n. Kiểm tra số dương n có phải là số hoàn thiện hay không?


    \textbf{\textit{Ghi chú}}: Một số nguyên dương $n$ được gọi là số hoàn thiện khi tổng các ước số thực sự của nó (các ước nhỏ hơn chính nó) bằng chính nó.

    \textit{Ví dụ minh họa}: Số 6 có ước số thực sự là $1, 2, 3$

    Tổng: $1 + 2 + 3 = 6$ $\rightarrow$ Là số hoàn thiện.
\end{problem}
\code{code/ch01/1-0-30.cpp}{c++}

\begin{problem}
    Cho số nguyên dương n. Kiểm tra số nguyên dương n có phải là số nguyên tố hay không?
\end{problem}
\code{code/ch01/1-0-31.cpp}{c++}

\begin{problem}
    Cho số nguyên dương n. Kiểm tra số nguyên dương n có phải là số chính phương hay không?

    \textbf{\textit{Ghi chú:}} Số $n$ là số chính phương nếu căn bậc hai của nó ($\sqrt{n}$) là một số nguyên.
\end{problem}
\code{code/ch01/1-0-32.cpp}{c++}

\begin{problem}
    Tính $S(n) = \sqrt{2 + \sqrt{2 + \sqrt{2+ \cdots + \sqrt{2 + \sqrt{2}}}}}$ có n dấu căn
\end{problem}
\code{code/ch01/1-0-33.cpp}{c++}

\begin{problem}
    Tính $S(n) = \sqrt{n + \sqrt{n - 1 + \sqrt{n - 2 + \cdots + \sqrt{2 + \sqrt{1}}}}}$ có n dấu căn
\end{problem}
\code{code/ch01/1-0-34.cpp}{c++}

\begin{problem}
    Tính $S(n) = \sqrt{1 + \sqrt{2 + \sqrt{3+ \cdots + \sqrt{n - 1 + \sqrt{n}}}}}$ có n dấu căn
\end{problem}
\code{code/ch01/1-0-35.cpp}{c++}

\begin{problem}
    Tính $S(n) = \sqrt{n! + \sqrt{(n-1)! + \sqrt{(n-2)! + \cdots + \sqrt{2! + \sqrt{1}}}}}$ có n dấu căn
\end{problem}
\code{code/ch01/1-0-36.cpp}{c++}

\begin{problem}
    Tính $S(n) = \sqrt[n]{n + \sqrt[n-1]{n-1 + \cdots + \sqrt[3]{3 + \sqrt{2}}}}$ có n-1 dấu căn
\end{problem}
\code{code/ch01/1-0-37.cpp}{c++}

\begin{problem}
    Tính $S(n) = \sqrt[n+1]{n + \sqrt[n]{n - 1 + \cdots + \sqrt[3]{2 + \sqrt{1}}}}$ có n dấu căn
\end{problem}
\code{code/ch01/1-0-38.cpp}{c++}

\begin{problem}
    Tính $S(n) = \sqrt[n+1]{n! + \sqrt[n]{(n - 1)! + \cdots + \sqrt[3]{2! + \sqrt{1!}}}}$ có n dấu căn
\end{problem}
\code{code/ch01/1-0-39.cpp}{c++}

\begin{problem}
    Tính $S(n,x) = \sqrt{x^n + \sqrt{x^{n - 1} + \sqrt{x^{n-2} + \cdots + \sqrt{x^2 + \sqrt{x}}}}}$ có n dấu căn
\end{problem}
\code{code/ch01/1-0-40.cpp}{c++}

\begin{problem}
    Tính $S(n) = \frac{1}{1 + (\frac{1}{1 +(\frac{1}{1 + (\frac{1}{1 + (\frac{1}{1 + 1})})})})}$ có n dấu phân số
\end{problem}
\code{code/ch01/1-0-41.cpp}{c++}

\begin{problem}
    Cho n là số nguyên dương. Hãy tìm giá trị nguyên dương k lớn nhất sao
cho $S(k) < n$. Trong đó chuỗi S(k) được định nghĩa như sau : $S(k) = 1 +
2 + 3 + … + k$.
\end{problem}
\code{code/ch01/1-0-42.cpp}{c++}

\begin{problem}
    Hãy đếm số lượng chữ số của số nguyên dương n.
\end{problem}
\code{code/ch01/1-0-43.cpp}{c++}

\begin{problem}
    Hãy tính tổng các chữ số của số nguyên dương n.
\end{problem}
\code{code/ch01/1-0-44.cpp}{c++}

\begin{problem}
    Hãy tính tích các chữ số của số nguyên dương n.
\end{problem}
\code{code/ch01/1-0-45.cpp}{c++}

\begin{problem}
    Hãy đếm số lượng chữ số lẻ của số nguyên dương n.
\end{problem}
\code{code/ch01/1-0-46.cpp}{c++}

\begin{problem}
    Hãy tính tổng các chữ số chẵn của số nguyên dương n.
\end{problem}
\code{code/ch01/1-0-47.cpp}{c++}

\begin{problem}
    Hãy tính tích các chữ số lẻ của số nguyên dương n.
\end{problem}
\code{code/ch01/1-0-48.cpp}{c++}

\begin{problem}
     Cho số nguyên dương n. Hãy tìm chữ số đầu tiên của n.
\end{problem}
\code{code/ch01/1-0-49.cpp}{c++}

\begin{problem}
    Hãy tìm chữ số đảo ngược của số nguyên dương n.
\end{problem}
\code{code/ch01/1-0-50.cpp}{c++}

\begin{problem}
    Tìm chữ số lớn nhất của số nguyên dương n.
\end{problem}
\begin{problem}
     Tìm chữ số nhỏ nhất của số nguyên dương n.
\end{problem}
\begin{problem}
    Hãy đếm số lượng chữ số lớn nhất của số nguyên dương n.
\end{problem}
\begin{problem}
     Hãy đếm số lượng chữ số nhỏ nhất của số nguyên dương n.
\end{problem}
\begin{problem}
     Hãy đêm số lượng chữ số đầu tiên của số nguyên dương n.
\end{problem}
\begin{problem}
    Hãy kiểm tra số nguyên dương n có toàn chữ số lẻ hay không?
\end{problem}
\begin{problem}
     Hãy kiểm tra số nguyên dương n có toàn chữ số chẵn hay không?
\end{problem}

\begin{problem}
    Hãy kiểm tra số nguyên dương n có phải số đối xứng hay không?
\end{problem}
\begin{problem}
    Hãy kiểm tra các chữ số của số nguyên dương n có tăng dần từ trái
sang phải hay không?
\end{problem}
\begin{problem}
    Hãy kiểm tra các chữ số của số nguyên dương n có giảm dần từ trái
sang phải hay không?
\end{problem}
\begin{problem}
    Cho hai số nguyên dương a và b. Hãy vẽ lưu đồ tìm ước chung lớn nhất
của hai giá trị này.
\end{problem}
\begin{problem}
    Cho hai số nguyên dương a và b. Hãy vẽ lưu đồ tìm bội chung nhỏ nhất
của hai giá trị này.
\end{problem}

\begin{problem}
    Giải phương trình $𝑎𝑥 + 𝑏 = 0$.
\end{problem}
\begin{problem}
    Giải phương trình $𝑎𝑥^2 + 𝑏𝑥 + 𝑐 = 0$.
\end{problem}
\begin{problem}
    Giải phương trình $𝑎𝑥^4 + 𝑏𝑥^2 + 𝑐 = 0$.
\end{problem}

\begin{problem}
    Tính $S(x, n) = 𝑥 − 𝑥^2 + 𝑥^3 + \cdots + (−1)^{n+1}𝑥^𝑛$
\end{problem}

\begin{problem}
    Tính $S(x, n) = −𝑥^2 + 𝑥^4 + \cdots + (−1)^n𝑥^{2𝑛}$
\end{problem}

\begin{problem}
    Tính $S(x, n) = 𝑥 − 𝑥^3 + 𝑥^5 + \cdots + (−1)^n𝑥^{2𝑛+1}$
\end{problem}

\begin{problem}
    Tính $S(n) = 1 - \frac{1}{1+2} + \frac{1}{1+2+3} + \cdots + (-1)^{n+1} \frac{1}{1+2+3+\cdots+n}$
\end{problem}

\begin{problem}
    Tính $S(x, n) = - x + \frac{x^2}{1+2} - \frac{x^3}{1+2+3} + \cdots + (-1)^{n} \frac{x^n}{1+2+3+\cdots+n}$
\end{problem}

\begin{problem}
    Tính $S(x, n) = -x + \frac{x^2}{2!} - \frac{x^3}{3!} + \cdots + (-1)^{n} \frac{x^n}{n!}$
\end{problem}

\begin{problem}
    Tính $S(x, n) = -1 + \frac{x^2}{2!} - \frac{x^4}{4!} + \cdots + (-1)^{n+1} \frac{x^{2n}}{2n!}$
\end{problem}

\begin{problem}
    Tính $S(x, n) = 1 - x + \frac{x^3}{3!} - \frac{x^5}{5!} + \cdots + (-1)^{n+1} \frac{x^{2n+1}}{{2n+1}!}$
\end{problem}

\begin{problem}
    Kiểm tra số nguyên 4 bytes có dạng $2^𝑘$ hay không?
\end{problem}

\begin{problem}
    Kiểm tra số nguyên 4 bytes có dạng $3^𝑘$ hay không?
\end{problem}

